
%\documentclass{report}
\documentclass[ams,openany,10pt,presentation,utf8]{common/themes/dossiers/dossiers}


\usepackage[object=vectorian]{pgfornament} %%  http://altermundus.com/pages/tkz/ornament/index.html



%%%%%%%%%%%%%%%%%%%%%%%%%%%%%%%%%%%%%%%%%%%%%%%%%%%%%%%%%%%%%%%% MARGES 
% method 2
\usepackage{geometry}
\geometry{hmargin=1.0cm,vmargin=1.5cm}
%\usepackage{geometry}
%\geometry{hmargin=1.5cm,vmargin=2.0cm}

%%%%%%%%%%%%%%%%%%%%%%%%%%%%%%%%%%%%%%%%%%%%%%%%%%%%%%%%%%%%%%%%


 

%%%%%%%%%%%%%%%%%%%%%%%%%%%%%%%%%%%%%%%%%%%%%%%%%%%%%%%%%%%%%%%% TIMING
\usepackage{tikz-timing}
\usetikzlibrary{calc}
\usetikzlibrary{positioning}


\usetikzlibrary{positioning,arrows.meta}
\usetikztiminglibrary[new={char=Q,reset char=R}]{counters}
\usetikztiminglibrary{nicetabs} 
\usetikzlibrary{shapes}


%%%%%%%%%%%%%%%%%%%%%%%%%%%%%%%%%%%%%%%%%%%%%%%%%%%%%%%%%%%%%%%%

%%%%%%%%%%%%%%%%%%%%%%%%%%%%%%%%%%%%%%%%%%%%%%%%%%%%%%%%%%%%%%%% HEADERS/FOOTERS
\usepackage{fancyhdr}
\pagestyle{fancy}
\usepackage{enumitem}



%\fancyhead[LE,RO]{\textsl{}}
%\fancyhead[LO,RE]{\textsl{}}

%\fancyhead[LE,RO]{\textsl{\rightmark}}
%\fancyhead[LO,RE]{\textsl{\leftmark}}
%\fancyhead[LE,RO]{\textsl{}}
\fancyfoot[C]{\thepage}


\fancyfoot[L]{MAOREDEV}
%\fancyfoot[C]{Electronique \& Communication}
%\fancyfoot[R]{\thepage}
\fancyfoot[R]{\today }
\renewcommand{\headrulewidth}{0.4pt}
%\renewcommand{\headrulewidth}{0.0pt}
\renewcommand{\footrulewidth}{0.4pt}

%%%%%%%%%%%%%%%%%%%%%%%%%%%%%%%%%%%%%%%%%%%%%%%%%%%%%%%%%%%%%%%%




%%%%%%%%%%%%%%%%%%%%%%%%%%%%%%%%%%%%%%%%%%%%%%%%%%%%%%%%%%%%%%%%

%%%%%%%%%%%%%%%%%%%%%%%%%%%%%%%%%%%%%%%%%%%%%%%%%%%%%%%%%%%%%%%%



\usepackage{amsfonts,amsmath,amssymb,amsthm}
\usepackage{aligned-overset}

\usepackage{mathptmx}
\usepackage{relsize}

\usepackage[T1]{fontenc}
\usepackage{graphicx}
\usepackage{ifthen}

\usepackage[Lenny]{fncychap}


\usepackage{listings,multicol}
\usepackage{color}

% print subsubsection
\setcounter{tocdepth}{5}
\setcounter{secnumdepth}{5}
% fin print subsubsection


\usepackage{minitoc}
\setcounter{minitocdepth}{3}
%\setcounter{minitocdepth}{2}
\setlength{\mtcindent}{24pt}
\renewcommand{\mtcfont}{\small\rm}
\renewcommand{\mtcSfont}{\small\bf}

\usepackage{hyperref}
\hypersetup{pdftex,colorlinks=true,allcolors=blue}
\usepackage{hypcap}


\usepackage{type1cm}
\usepackage{lettrine}
\usepackage{scrextend}
\usepackage{blindtext}
%\usepackage[svgnames, x11names]{xcolor}



\usepackage[utf8]{inputenc}
%\usepackage{lmodern,textcomp}

\usepackage[french]{babel}
\mtcselectlanguage{french}

\usepackage{multido}
\RequirePackage{framed}
\usepackage{fancyhdr}
\usepackage{fancybox}
\usepackage{fancybox,xcolor}
\usepackage{blindtext}
\usepackage{xcolor, colortbl} %kleurtjes voor table

\usepackage{hhline}
\usepackage{makecell} % https://tex.stackexchange.com/questions/2441/how-to-add-a-forced-line-break-inside-a-table-cell



\usepackage{answers}

%\usepackage{mdframed}

\usepackage[framemethod=TikZ]{mdframed}

\usepackage{pstricks,pst-grad}

% ntheorem
\makeatletter
\let\th@plain\relax
\makeatother
%\usepackage[thref]{ntheorem}
%\usepackage[amsmath,thref]{ntheorem}
\usepackage{empheq}
\usepackage{thmtools}
\usepackage{subfig}
\usepackage{hyperref}%
\usepackage{systeme,mathtools}



%%%%%%%%%%%%%%%%%%%%%%%%%%%%%%%%%%%%%%%%%%%%%%%%%%%%%%%%%%%%%%%% ARRAY 
\usepackage{array,multirow,makecell}
\renewcommand{\arraystretch}{1.2}
%\setcellgapes{1pt}
%\makegapedcells
%\newcolumntype{R}[1]{>{\raggedleft\arraybackslash }b{#1}}
%\newcolumntype{L}[1]{>{\raggedright\arraybackslash }b{#1}}
%\newcolumntype{C}[1]{>{\centering\arraybackslash }b{#1}}
%%%%%%%%%%%%%%%%%%%%%%%%%%%%%%%%%%%%%%%%%%%%%%%%%%%%%%%%%%%%%%%%

%%%%%%%%%%%%%%%%%%%%%%%%%%%%%%%%%%%%%%%%%%%%%%%%%%%%%%%%%%%%%%%% PERSONNAL
\usepackage{common/lesson/headers}
\usepackage{common/lesson/blocks}
\usepackage{common/lesson/theorems}
%\usepackage{common/lesson/tests}
\usepackage{common/lesson/ccf}
\usepackage{common/lesson/pretty}

\usepackage{common/images/nodes}
\usepackage{common/lesson/tables}

\definecolor{Gray}{gray}{0.9}

%%%%%%%%%%%%%%%%%%%%%%%%%%%%%%%%%%%%%%%%%%%%%%%%%%%%%%%%%%%%%%%%

% moi-même
\renewcommand{\thesection}{\arabic{section}} % Remove chapter number to sections

\usepackage{chngcntr}
\usepackage{xhfill}


% Circle
\usepackage{pict2e,picture}
\newsavebox\CBox
\newlength\CLength
\def\competence#1{\sbox\CBox{#1}%
  \ifdim\wd\CBox>\ht\CBox \CLength=\wd\CBox\else\CLength=\ht\CBox\fi
    \makebox[1.1\CLength]{\makebox(0,1.1\CLength){\put(0,0){\circle{1.0\CLength}}\put(0,0){\circle{0.9\CLength}}}%
    \makebox(0,1.1\CLength){\put(-.3\wd\CBox,0){\tiny\textbf{#1}}}}}
% End circle


%----- Ensembles : entiers, reels, complexes -----
\newcommand{\Nn}{\mathbb{N}} \newcommand{\N}{\mathbb{N}}
\newcommand{\Zz}{\mathbb{Z}} \newcommand{\Z}{\mathbb{Z}}
\newcommand{\Qq}{\mathbb{Q}} \newcommand{\Q}{\mathbb{Q}}
\newcommand{\Rr}{\mathbb{R}} \newcommand{\R}{\mathbb{R}}
%\newcommand{\Cc}{\mathbb{C}} \newcommand{\C}{\mathbb{C}}
\newcommand{\Kk}{\mathbb{K}} \newcommand{\K}{\mathbb{K}}


\newcommand{\sk}{\vspace{0.2 cm}}
\newcommand{\A}[1]{{$\backslash${\tt #1}}}
\newcommand{\nsp}{\mbox{\hspace{-1 cm}}}

\title{\LARGE \textbf{BTS Systèmes Numériques} \\ \normalsize (Option \'Electronique \& Communication) \\ \Large CCF 2020-2022 \\ Sujet X}
%\author{ABDULLATIF Mouhamadi}
\author{CFAI Poitou-Charentes}
\date{\today}



%\def\PRINTMODE{PROFESSEUR}
\def\PRINTMODE{ELEVE}


\newcommand{\TEACHERBLOC}{\ifthenelse{\equal{\PRINTMODE}{PROFESSEUR}}}



%%%%%%%%%%%%%%%%%%%%%%%%%%%%%%%%%%%%%%%%%%%%%%%%%%%%%%%%% remove chapter
\usepackage{titlesec}
%\titleformat{\chapter}[display]{\normalfont\bfseries}{}{0pt}{\Huge}
  %{\normalfont\bfseries}{}{0pt}{\Huge}
\usepackage{ulem}
\usepackage[subfigure]{tocloft}

\usepackage[final]{pdfpages}
\usepackage{pdflscape}






\AtBeginDocument{
    \renewcommand{\labelitemi}{\textbullet}
    \renewcommand{\labelitemi}{$\circ$}
    \renewcommand{\labelitemii}{$\circ$}
    \renewcommand{\labelitemiii}{$\circ$}
    \renewcommand{\labelitemiv}{$\circ$}
}



%%%%%%%%%%%%%%%%%%%%%%%%%%%%%%%%%%%%%%%%%%%%%%%%%%%%%%%%%%%%%%%% IMAGE FRAME 
\usepackage{efbox,graphicx}
%_______________________________________________________________

\usepackage{csvsimple}

\usepackage{longtable}


\newcommand{\redt}[1]{\textcolor{general@best@maroon}{\textbf{#1}}}

%yay -S texlive-tikz-uml
\usepackage{tikz-uml}
%\usepackage{umlpackage}
\usepackage{uml}






\newcommand{\horizb}{%
  \nointerlineskip \vspace{.5\baselineskip}\hspace{\fill}
  {\color{general@best@orange}
    \resizebox{0.75\linewidth}{2ex}
    {{%
    {\begin{tikzpicture}
    \node  (C) at (0,0) {};
    \node (D) at (9,0) {};
    \path (C) to [ornament=89] (D);
    \end{tikzpicture}}}}}%
    \hspace{\fill}
    \par\nointerlineskip \vspace{.5\baselineskip}
  }


%----------------------------------------------------------------------
\def\signalracolor{general@best@purple}
\newcommand{\signalra}[1]{
  \begin{tikzpicture} [every node/.style={signal, draw, text=orange, signal to=nowhere, inner sep=1.1pt}]
    %\node[white] (FAKE) at (0,.25) {};
    \node[white, fill=\signalracolor, signal to=east, signal from=nowhere] at (0,0) {\ttfamily\normalsize#1};
  \end{tikzpicture}
}


%----------------------------------------------------------------------
%.25-.7

\def\sourcelinkacolor{general@best@teal}
\newcommand{\sourcelinka}[1]{%
  \TEACHERBLOC{
    \begin{longtable}{cc}
      \arrayrulecolor{\sourcelinkacolor}
      %\hline
      \hline 
      \rowcolor{white}\multicolumn{2}{|c|}{ \textcolor{\sourcelinkacolor}{\textsc{sources}} }  \\
      \rowcolor{\sourcelinkacolor} \bfseries \textcolor{white}{thème} & \bfseries \textcolor{white}{lien}\\\hline
      %\rowcolor{\sourcelinkacolor}\multicolumn{2}{C{\textwidth}}{ \textbf{sources} } \\
      %\rowcolor{\sourcelinkacolor}\multicolumn{1}{C{.25\textwidth}}{ \textbf{titre} } & \multicolumn{1}{C{.7\textwidth}|}{ \textbf{lien} } \\
      #1
      %\hline
    \end{longtable}
  }{}
}
\newcommand{\slar}[2]{%
%#1&#2&#3\\\hline
\multicolumn{1}{L{.25\textwidth}}{ \cellcolor{\sourcelinkacolor!16}\textit{#1} } & \multicolumn{1}{L{.7\textwidth}}{ \footnotesize\url{#2} } \\\hline
%\cellcolor{general@best@aqua}{#1}&\footnotesize\url{#2}\\%\hline
}


% ----------------------------------------------------------------------

\def\articlesacolor{general@best@maroon}
\newcommand{\articlesa}[1]{%
    \begin{longtable}{cccccc}
      \arrayrulecolor{\articlesacolor}
      \hline 
      \rowcolor{white}\multicolumn{6}{|c|}{ \textcolor{\articlesacolor}{\textsc{catalogue}} }  \\
      \rowcolor{\articlesacolor} \bfseries \textcolor{white}{$\boldsymbol{\#}$}&
      \textcolor{white}{\bfseries image}&
      \textcolor{white}{\bfseries article}&
      \textcolor{white}{\bfseries commentaires}&
      \textcolor{white}{\bfseries prix}&
      \textcolor{white}{\bfseries lien}\\\hline
      #1
    \end{longtable}
}
\newcommand{\artr}[6]{%
\multicolumn{1}{|C{.05\textwidth}}{#1} &
\multicolumn{1}{|C{.20\textwidth}}{\includegraphics[width=1\linewidth]{#2}} &
\multicolumn{1}{|L{.270\textwidth}}{\footnotesize\bfseries #3}&
\multicolumn{1}{|L{.220\textwidth}}{\footnotesize\bfseries #4}&
\multicolumn{1}{|C{.075\textwidth}}{\textcolor{red}{\bfseries #5€}}&
\multicolumn{1}{|C{.05\textwidth}|}{\large\href{#6}{aller}}\\\hline
}

%    $1$&\includegraphics[width=\linewidth]{img1}&\footnotesize\bfseries Wangkangyi Moniteur portable 15" TFT-LCD USB écran tactile LCD PC Moniteur de caisse utilisé pour Kiosk Restaurant Bar&167&\large\href{https://www.amazon.fr/dp/B08GCGV99B/}{aller}

%|C{.05\textwidth}|L{.3\textwidth}|L{.393\textwidth}|C{0.075\textwidth}|C{.05\textwidth}|



% ----------------------------------------------------------------------
\def\pricingacolor{general@best@orange}
\newcommand{\pricinga}[1]{%
    \begin{longtable}{ccccc}
      \arrayrulecolor{\pricingacolor}
      \hline 
      \rowcolor{white}\multicolumn{5}{|c|}{ \textcolor{\pricingacolor}{\textsc{pricing}} }  \\
      \rowcolor{\pricingacolor} \bfseries \textcolor{white}{$\boldsymbol{\#}$}&
      \textcolor{white}{\bfseries service}&
      \textcolor{white}{\bfseries fixe (h)}&
      \textcolor{white}{\bfseries /mois (€)}&
      \textcolor{white}{\bfseries commentaires}\\\hline
      #1
    \end{longtable}
}
\newcommand{\prr}[5]{%
\multicolumn{1}{|C{.05\textwidth}}{#1} &
\multicolumn{1}{|L{.25\textwidth}}{\bfseries #2} &
\multicolumn{1}{|C{.10\textwidth}}{\bfseries \textcolor{red}{#3h}} &
\multicolumn{1}{|C{.10\textwidth}}{\bfseries \textcolor{red}{#4€}} &
\multicolumn{1}{|L{.38\textwidth}}{\footnotesize\bfseries #5}\\\hline
}



% ----------------------------------------------------------------------



\def\tutorialacolor{general@best@navy}
\newcommand{\tutoriala}[2]{%
\begin{longtable}{cc}
  \arrayrulecolor{\tutorialacolor}
  \hline
  \rowcolor{white}\multicolumn{2}{|c|}{ \textcolor{\tutorialacolor}{\textsc{#1}} }  \\
  \hline

  #2
\end{longtable}
}

\newcommand{\tutorialapic}[2]{%
\rowcolor{\tutorialacolor!16}\multicolumn{2}{c}{ \raisebox{-\totalheight}{
  \begin{tikzpicture}
      \node[inner sep=0pt] (img) at (0,0) {\includegraphics[width=0.95\linewidth]{#1}};

      #2
    \end{tikzpicture}
  }}\\
  \rowcolor{\tutorialacolor!16}\multicolumn{2}{c}{~} \\\hline
}
\newcommand{\tutorialatag}[2]{%
  \node[white, chamfered rectangle,double=\tutorialacolor, draw, fill=\tutorialacolor, very thick] at #1 {\bfseries #2};
}
\newcommand{\tutorialarow}[2]{%
  \multicolumn{1}{C{.05\textwidth}|}{ \footnotesize\charmed{#1} } & \multicolumn{1}{L{.9\textwidth}}{ #2 } \\\hline
}

% ----------------------------------------------------------------------

\newcommand*\circled[1]{\tikz[baseline=(char.base)]{%
            \node[white,shape=circle,fill=general@best@navy!75,draw,inner sep=2.0pt, very thick] (char) {\bfseries #1};}}
\newcommand*\stared[1]{\tikz[baseline=(char.base)]{%
            \node[white,shape=star, star points=8,fill=general@best@navy!75,draw,inner sep=2.0pt, very thick] (char) {\bfseries #1};}}
\newcommand*\charmed[1]{\tikz[baseline=(char.base)]{%
            \node[white, chamfered rectangle,double=\tutorialacolor, draw, fill=\tutorialacolor, very thick] (char) {\bfseries #1};}}




% https://tex.stackexchange.com/questions/22406/how-to-make-a-textbox-with-this-tikz-code
\colorlet{bancolor}{general@best@navy}
\def\bancolor{mycolor}
\newenvironment{objectifa}[3][]{%
    \begin{tikzpicture}[#1]%
        \def\myboxname{\textcolor{white}{\textsc{#3}}}%
        \node [bancolor,draw,inner sep=1.5ex,text width=#2]% good options: minimum height, minimum width
            (BOXCONTENT) \bgroup\rule{0pt}{0ex}\ignorespaces
}{%
        \egroup;
        \node [right,inner xsep=1em,fill=bancolor!75,outer sep=0pt,text height=2ex,text depth=.5ex] (BOXNAME) 
            at ([shift={(-1em,0pt)}]BOXCONTENT.north west) {\myboxname};
        \fill[bancolor] (BOXNAME.north east) -- +(-1em,1em) -- +(-1em,0) -- cycle;
        \fill[bancolor] (BOXNAME.south west) -- +(1em,-1em) -- +(1em,0) -- cycle;
    \end{tikzpicture}
}

% chamfered rectangle
% ----------------------------------------------------------------------

\newcommand{\bubblea}[1]{
  \begin{tikzpicture}
    \node[fill=general@best@teal,rounded corners]
    {\textcolor{white}{\bfseries#1}};
    \end{tikzpicture}
}
\newcommand{\bubbleb}[1]{
  \begin{tikzpicture}
    \node[general@best@olive, draw,inner sep=1.0pt]{#1};
    \end{tikzpicture}
}
\newcommand{\bubblec}[1]{
  \begin{tikzpicture}
    \node[general@best@maroon, draw, inner sep=1.0pt]{#1};
    \end{tikzpicture}
}
\newcommand{\bubblefa}[1]{
  \begin{tikzpicture}
    \node[fill=general@best@red,rounded corners]
    {\textcolor{white}{\bfseries#1}};
    \end{tikzpicture}
}
\newcommand{\bubblefb}[1]{
  \begin{tikzpicture}
    \node[fill=general@best@olive,rounded corners]
    {\textcolor{white}{\bfseries#1}};
    \end{tikzpicture}
}